\chapter{Web apps}\label{ch:web}

Three helpers run the web apps on the Mac or a home server, defined in
\file{zsh/.config/zsh/50-web.zsh}:

\begin{keys}{3.6cm}
\head{Helper}{For}
\cmd{app} & Start an app in its own tmux window, from a short list of your projects\\
\cmd{tun} & See and change what the Cloudflare Tunnel serves, or share a port for a moment\\
\cmd{pml} \cmd{pmlog} \cmd{pmr} \cmd{pms} & PM2 shortcuts\\
\end{keys}

\section{The app list}

\file{install.sh} creates \file{~/.config/apps}. One line per app: a name, its folder, and
optionally the command that starts it.

\begin{plain}
# name        folder                      [command]
molly         ~/apps/molly-dashboard
gcs           ~/apps/nextgen-gcs
netra         ~/apps/netra                scripts/local.sh up
database      ~/apps/database
\end{plain}

Without a command, the folder decides how the app starts:

\begin{keys}{5.2cm}
\head{The folder has}{\texttt{app} runs}
\file{ecosystem.config.js} (or \file{.cjs}, \file{.json}) & \cmd{pm2 start <file>}, then follows \cmd{pm2 logs}\\
\file{compose.yaml} or \file{docker-compose.yml} & \cmd{docker compose up -d}, then follows the logs\\
\file{package.json} & \cmd{pnpm dev}, \cmd{bun run dev}, \cmd{yarn dev} or \cmd{npm run dev}, picked from the lockfile\\
none of these & Only opens a shell in the folder\\
\end{keys}

\section{The app command}

\begin{keys}{3.6cm}
\head{Command}{Action}
\cmd{app} or \key{C-a a} & Pick an app: a new tmux window named after it opens in its folder and starts it\\
\cmd{app <name>} & The same without the picker\\
\cmd{app ls} & Every app with how it starts (\texttt{pm2}, \texttt{compose}, \texttt{node}, \texttt{cmd}) and its folder\\
\cmd{app up <name>} & Start in the background (PM2 and compose)\\
\cmd{app down <name>} & \cmd{pm2 stop} or \cmd{docker compose down}\\
\cmd{app logs <name>} & PM2 or compose logs\\
\cmd{app cd <name>} & Change into the folder\\
\end{keys}

\screenshot[0.7]{app-picker}{\key{C-a a}: the app picker, with the folder of the selected app on the right.}

A typical morning:

\begin{shell}
t                 # tmux session main
app molly         # window "molly": pnpm dev
app database      # window "database": docker compose up -d, then the logs
tun               # are the public hostnames up?
\end{shell}

\section{Cloudflare Tunnel}

\cmd{tun} works on a locally managed tunnel: one \file{config.yml} with ingress rules, run by
\texttt{cloudflared} as a service. Tunnels managed in the Cloudflare dashboard (token based) have no
local config; change those in the dashboard.

\subsection{One-time setup}

\begin{shell}
cloudflared tunnel login          # browser login, writes ~/.cloudflared/cert.pem
cloudflared tunnel create mac     # writes ~/.cloudflared/<tunnel-id>.json
\end{shell}

Then write \file{~/.cloudflared/config.yml}:

\begin{plain}
tunnel: <tunnel-id>
credentials-file: /Users/<you>/.cloudflared/<tunnel-id>.json

ingress:
  - hostname: molly.example.com
    service: http://localhost:3000
  - service: http_status:404      # must stay last: tun adds routes above it
\end{plain}

\begin{shell}
cloudflared tunnel route dns mac molly.example.com
brew services start cloudflared   # starts at login, reads ~/.cloudflared/config.yml
\end{shell}

\begin{warning}
Start the service as your user, not with \cmd{sudo}: a root service reads its config from
\file{/etc/cloudflared}, not from your home folder.
\end{warning}

\subsection{Commands}

\begin{keys}{4.2cm}
\head{Command}{Action}
\cmd{tun} & Running or not, the tunnel name, and every hostname \(\rightarrow\) service route\\
\cmd{tun add <host> <port>} & Add a route (a bare port becomes \texttt{http://localhost:<port>}), create its DNS record, restart\\
\cmd{tun rm <host>} & Remove a route and restart (delete the DNS record in the dashboard)\\
\cmd{tun test <host>} & Which rule a request for that hostname would hit\\
\cmd{tun share <port>} & Temporary \texttt{*.trycloudflare.com} URL for a local port; \key{C-c} ends it\\
\cmd{tun log} & Follow the cloudflared log\\
\cmd{tun restart} & Restart the service\\
\cmd{tun edit} & Open \file{config.yml} in Neovim\\
\end{keys}

\screenshot[0.7]{tun}{\cmd{tun}: tunnel status and the hostname routes.}

Every change copies the config to \file{config.yml.bak} first and runs
\cmd{cloudflared tunnel ingress validate} afterwards; an invalid result is rolled back.

Settings, in \file{~/.zshrc.local}:

\begin{keys}{4.6cm}
\head{Variable}{Default}
\texttt{CLOUDFLARED\_CONFIG} & \file{~/.cloudflared/config.yml}\\
\texttt{CLOUDFLARED\_RESTART} & \cmd{brew services restart cloudflared}\\
\end{keys}

With \cmd{cloudflared service install} instead of Homebrew:

\begin{shell}
export CLOUDFLARED_RESTART='launchctl kickstart -k gui/$(id -u)/com.cloudflare.cloudflared'
\end{shell}

\section{PM2 and Docker}

\begin{keys}{3.8cm}
\head{Alias}{Command}
\cmd{pml} & \cmd{pm2 ls}\\
\cmd{pmlog} & \cmd{pm2 logs --lines 100}\\
\cmd{pmr <name>} & \cmd{pm2 restart <name>}\\
\cmd{pms} & \cmd{pm2 save} (the list that \cmd{pm2 resurrect} restores)\\
\cmd{dps}, \cmd{dpsa}, \cmd{di} & Containers, all containers, images\\
\cmd{dcu}, \cmd{dcd}, \cmd{dcl}, \cmd{dcb} & \cmd{docker compose} up -d, down, logs -f, build\\
\cmd{dsh [container]} & Shell into a container (zsh, bash or sh), fuzzy pick without a name\\
\cmd{dlog}, \cmd{dstat} & Follow logs, live CPU/RAM per container\\
\cmd{dstop}, \cmd{drm} & Stop / remove containers (multi-select with \key{Tab})\\
\cmd{dclean} & Prune stopped containers, networks, dangling images and volumes\\
\cmd{lzd} or \key{C-a D} & lazydocker\\
\end{keys}

To start the saved PM2 apps at login: \cmd{pm2 startup} once, then \cmd{pm2 save} whenever the list
changes.
